\chapter{Algoritmos de Simulación}
\label{ane_00_00_anexo}
Este anexo presenta los algoritmos empleados para generar los datos del estudio de simulación del Capítulo~\ref{03_00_00_mezclas_ddp_series_tiempo_funcionales}. Cada algoritmo especifica su mecanismo generador, el rol de sus parámetros y los valores numéricos con que se emplea, los cuales se recogen en un cuadro al cierre de cada uno.\\

Los algoritmos se agrupan según el nivel en que se especifica la dinámica del proceso:
\begin{itemize}
    \item Métodos basados en series de tiempo clásicas.
    \item Métodos basados en procesos funcionales.
    \item Métodos basados en coeficientes de bases.
\end{itemize}

\section{Métodos basados en Series de Tiempo Clásicas}
\label{ane_00_01_metodos_series_tiempo}
Este primer enfoque genera las curvas a partir de una serie de tiempo escalar ``clásica''. La idea es simular una única trayectoria muy larga de un proceso univariado discreto $\{Z_s\}$ y luego segmentarla en bloques consecutivos de igual longitud, de modo que cada bloque se interprete como una curva. Esto reproduce la manera en que suelen construirse series de tiempo funcionales a partir de registros de alta frecuencia, donde el registro se divide en sus unidades naturales de observación (por ejemplo, el día) y cada unidad se trata como una realización funcional.\\

Los escenarios de esta sección comparten el mecanismo de segmentación y difieren únicamente en el proceso escalar generador. Sea $L$ el número de puntos de la grilla, que en este enfoque coincide con la longitud de cada bloque. Cada algoritmo simula una trayectoria $Z_1, \dots, Z_{TL}$ y la divide en $T$ bloques contiguos y sin traslape, de modo que la curva $t$ queda formada por $Z_{(t-1)L+1}, \dots, Z_{tL}$ sobre la grilla $\{\tau_l\}_{l=1}^{L}$. Sobre las curvas así obtenidas se aplica el esquema de observación común.\\

Los escenarios definidos son los siguientes:
\begin{itemize}
    \item El Algoritmo~A-1 induce dependencia de largo alcance entre curvas mediante un proceso ARFIMA.
    \item El Algoritmo~A-2 induce una ley condicional de mezcla mediante un proceso autorregresivo con cambio de régimen markoviano latente.
    \item El Algoritmo~A-3 induce heterocedasticidad condicional persistente mediante un proceso autorregresivo con innovaciones GARCH.
\end{itemize}

\subsection{Algoritmo A-1: Proceso ARFIMA de memoria larga}
\label{ane_00_01_01_alg_a1}
Este algoritmo genera curvas a partir de un proceso ARFIMA$(0,d,0)$ \cite{ane_00_01_granger1980introduction, ane_00_01_hosking1981fractional}, con el fin de que la dependencia entre curvas decaiga lentamente con la distancia. El proceso se define como
\begin{equation}\label{eq:ane_algA1_arfima}
    (1-B)^d\, Z_s = \zeta_s,
    \qquad
    \zeta_s \overset{\mathrm{iid}}{\sim} \mathcal{N}(0,\, \sigma_\zeta^2),
\end{equation}
donde $B$ es el operador de rezago y $(1-B)^d$ es el operador de diferenciación fraccionaria. El parámetro $d \in (0, 1/2)$ controla la memoria del proceso, garantiza estacionariedad y hace que la autocorrelación decaiga de forma hiperbólica, $\rho(k) \sim \frac{\Gamma(1-d)}{\Gamma(d)}\, k^{2d-1}$, y no geométrica como en un proceso ARMA. La autocorrelación al rezago $L$ gobierna la dependencia entre puntos homólogos de curvas consecutivas, y crece con $d$. Se calcula de forma recursiva mediante
\begin{equation}\label{eq:ane_algA1_acf}
    \rho(0) = 1,
    \qquad
    \rho(k) = \rho(k-1)\,\frac{k-1+d}{k-d}, \quad k \geq 1,
\end{equation}
y la varianza de la innovación se despeja como $\sigma_\zeta^2 = \sigma_Z^2\,\Gamma(1-d)^2/\Gamma(1-2d)$, de modo que $\mathrm{Var}(Z_s) = \sigma_Z^2$ para todo $d$.\\

La generación procede de la siguiente forma:
\begin{enumerate}
    \item Se fijan $d$ y $\sigma_Z$, y se calculan las autocovarianzas $\gamma(k) = \sigma_Z^2\,\rho(k)$, $k = 0, \dots, TL-1$, mediante \eqref{eq:ane_algA1_acf}.

    \item Se simula la trayectoria $Z_1, \dots, Z_{TL}$ de forma exacta mediante el método de Davies--Harte \cite{ane_00_01_davies1987tests}, que incrusta la matriz de covarianza en una matriz circulante y utiliza la transformada rápida de Fourier.
\end{enumerate}

\begin{table}[H]
\centering
\caption{Parámetros del Algoritmo~A-1.}
\label{tab:ane_algA1}
\begin{tabular}{llc}
\hline
Símbolo & Descripción & Valor \\
\hline
$d$          & parámetro de memoria larga                & $0.4$ \\
$\sigma_Z$   & desviación estándar marginal de $Z_s$     & $0.5$ \\
\hline
\end{tabular}
\end{table}

\subsection{Algoritmo A-2: Proceso autorregresivo con cambio de régimen markoviano}
\label{ane_00_01_02_alg_a2}
Este algoritmo genera curvas a partir de un proceso autorregresivo con cambio de régimen markoviano \cite{ane_00_01_hamilton1989new}, con el fin de inducir una ley condicional de mezcla cuyo régimen persiste a lo largo de varias curvas. El proceso se define como
\begin{equation}\label{eq:ane_algA2_ms}
    Z_s = m_{S_s} + u_s,
    \qquad
    u_s = \phi\, u_{s-1} + \sigma_{S_s}\, \varepsilon_s,
    \qquad
    \varepsilon_s \overset{\mathrm{iid}}{\sim} \mathcal{N}(0,1),
\end{equation}
donde $\{S_s\}$ es una cadena de Markov independiente de $\{\varepsilon_s\}$, con espacio de estados $\{1,2\}$ y matriz de transición simétrica,
\begin{equation}\label{eq:ane_algA2_transicion}
    \mathbf{P} =
    \begin{pmatrix} p & 1-p \\ 1-p & p \end{pmatrix}.
\end{equation}
Los regímenes se distinguen por su nivel, $m_1 = -m$ y $m_2 = +m$, y por su escala, $\sigma_1 < \sigma_2$, mientras que $\phi$, con $|\phi|<1$, es común a ambos y determina la suavidad de la trayectoria dentro de cada régimen. Cada régimen dura en promedio $1/(1-p)$ observaciones, es decir, $1/\bigl((1-p)L\bigr)$ curvas; por ello $p$ se toma cercano a uno, pues la persistencia de la cadena a escala de curva es lo que induce dependencia entre curvas consecutivas.\\

La generación procede de la siguiente forma:
\begin{enumerate}
    \item Se simula la cadena $S_1, \dots, S_{TL}$, extrayendo $S_1$ de su distribución estacionaria, $(1/2, 1/2)$, y cada estado posterior desde la fila de $\mathbf{P}$ correspondiente al estado anterior.

    \item Se inicializa $u_0 = 0$ y se itera \eqref{eq:ane_algA2_ms}, obteniendo $Z_1, \dots, Z_{TL}$. La cadena latente $\{S_s\}$ no se entrega a los métodos.
\end{enumerate}
\begin{table}[H]
\centering
\caption{Parámetros del Algoritmo~A-2.}
\label{tab:ane_algA2}
\begin{tabular}{llc}
\hline
Símbolo & Descripción & Valor \\
\hline
$p$                 & probabilidad de permanencia en el régimen     & $0.998$ \\
$m$                 & desplazamiento de nivel, $m_1=-m$, $m_2=+m$   & $1.5$ \\
$\sigma_1,\sigma_2$ & escala de la innovación en cada régimen       & $0.5$, $1.0$ \\
$\phi$              & coeficiente autorregresivo común              & $0.5$ \\
\hline
\end{tabular}
\end{table}

\subsection{Algoritmo A-3: Proceso AR-GARCH}
\label{ane_00_01_03_alg_a3}
Este algoritmo genera curvas a partir de un proceso autorregresivo con innovaciones GARCH$(1,1)$ \cite{ane_00_01_bollerslev1986generalized}, con el fin de inducir heterocedasticidad condicional persistente: la dependencia entre curvas consecutivas opera sobre la escala de la serie y no sobre su nivel. El proceso se define como
\begin{equation}\label{eq:ane_algA3_argarch}
    Z_s = \phi\, Z_{s-1} + a_s,
    \qquad
    a_s = \sqrt{h_s}\, z_s,
    \qquad
    z_s \overset{\mathrm{iid}}{\sim} \mathcal{N}(0,1),
\end{equation}
con varianza condicional
\begin{equation}\label{eq:ane_algA3_recursion}
    h_s = \alpha_0 + \alpha_1\, a_{s-1}^2 + \beta_1\, h_{s-1},
\end{equation}
donde $\alpha_1$ gobierna la sensibilidad de la varianza futura a la innovación más reciente y $\beta_1$ la persistencia de la varianza condicional pasada. La estacionariedad de segundo orden requiere $\alpha_1 + \beta_1 < 1$, y esta persistencia total determina cuántas curvas se extiende la agrupación de volatilidad. El nivel base se despeja de la varianza incondicional de la innovación como $\alpha_0 = \sigma_a^2\,(1-\alpha_1-\beta_1)$, y $\phi$, con $|\phi|<1$, determina la suavidad de la trayectoria dentro de cada curva.\\

La generación procede de la siguiente forma:
\begin{enumerate}
    \item Se fijan $\phi$, $\alpha_1$, $\beta_1$ y $\sigma_a^2$, y se calcula $\alpha_0$.

    \item Se inicializan $h_1 = \sigma_a^2$ y $Z_0 = 0$, y se itera \eqref{eq:ane_algA3_argarch}--\eqref{eq:ane_algA3_recursion}, obteniendo $Z_1, \dots, Z_{TL}$.
\end{enumerate}

\begin{table}[H]
\centering
\caption{Parámetros del Algoritmo~A-3.}
\label{tab:ane_algA3}
\begin{tabular}{llc}
\hline
Símbolo & Descripción & Valor \\
\hline
$\phi$          & coeficiente autorregresivo                          & $0.5$ \\
$\alpha_1$      & sensibilidad a la innovación reciente               & $0.02$ \\
$\beta_1$       & persistencia de la varianza condicional pasada      & $0.975$ \\
$\sigma_a^2$    & varianza incondicional de la innovación             & $1.0$ \\
\hline
\end{tabular}
\end{table}

\section{Métodos basados en Procesos Funcionales}
\label{ane_00_02_metodos_procesos_funcionales}

Este segundo enfoque genera directamente las curvas funcionales y permite que la relación entre observaciones consecutivas dependa de distintos mecanismos de generación. A diferencia de los métodos basados en series de tiempo clásicas, la dinámica se especifica sobre la curva completa, de modo que cada período puede ser generado mediante un mecanismo funcional diferente. En particular, los escenarios consideran una estructura de mezcla en la que una variable latente determina el mecanismo que genera la curva siguiente.\\

Sea $Z_t \in \{1,\ldots,K\}$ el mecanismo asociado al período $t$. Condicionalmente a $Z_t=k$, la curva se genera mediante:

\begin{equation}
\label{eq:ane_algB_general}
X_t(\tau)
=
f_k\!\left(R_{t-1,1},\ldots,R_{t-1,p}\right)(\tau)
+
\varepsilon_t(\tau),
\end{equation}

donde $f_k$ representa el mecanismo funcional correspondiente al estado $k$, $R_{t-1,j}$ denota el $j$-ésimo rasgo obtenido de la curva anterior y $\varepsilon_t(\tau)$ corresponde al término de error funcional.\\

Los escenarios se organizan de forma progresiva según la complejidad de esta estructura:

\begin{itemize}
    \item El Algoritmo~B-1 considera varios mecanismos funcionales con relaciones lineales, de modo que la heterogeneidad se encuentra en los parámetros de generación de cada mecanismo.

    \item El Algoritmo~B-2 incorpora mecanismos con relaciones no lineales, permitiendo que distintos mecanismos representen transformaciones lineales y no lineales de los rasgos funcionales de la curva anterior.

    \item El Algoritmo~B-3 combina mecanismos no lineales con una asignación dependiente de múltiples características de la curva anterior, de modo que la información disponible en $X_{t-1}$ determina probabilísticamente qué mecanismo genera $X_t$.
\end{itemize}

\subsection{Algoritmo B-1: Mezcla de mecanismos funcionales lineales}
\label{ane_00_02_01_alg_b1}

Este algoritmo genera curvas mediante una mezcla de mecanismos funcionales lineales. El objetivo es introducir heterogeneidad en la relación entre curvas consecutivas sin incorporar todavía efectos no lineales. Para ello, se consideran $K$ mecanismos de generación, cada uno caracterizado por una función media y un conjunto de coeficientes funcionales.\\

Sea $Z_t \in \{1,\ldots,K\}$ la variable latente que identifica el mecanismo utilizado en el período $t$. Condicionalmente a $Z_t=k$, la curva se genera como

\begin{equation}
\label{eq:ane_algB1}
X_t(\tau)
=
\mu_k(\tau)
+
\sum_{j=1}^{p}
\beta_{kj}(\tau)R_{t-1,j}
+
\varepsilon_t(\tau),
\end{equation}

donde $\mu_k(\tau)$ representa la función media asociada al mecanismo $k$, $\beta_{kj}(\tau)$ es el efecto funcional del $j$-ésimo rasgo bajo dicho mecanismo y $R_{t-1,j}$ corresponde a una característica de la curva anterior. En este escenario se consideran $p=3$ rasgos funcionales, de manera que cada curva futura depende simultáneamente de tres características de $X_{t-1}$.\\

Los rasgos se definen mediante funcionales lineales de la curva anterior,

\begin{equation}
\label{eq:ane_algB1_rasgos}
R_{t-1,j}
=
\int_0^1 X_{t-1}(s)w_j(s)\,ds,
\qquad j=1,2,3,
\end{equation}

donde $w_j(s)$ son funciones de ponderación prefijadas. De esta forma, cada rasgo resume información de una región o característica particular de la curva anterior.\\

La variable $Z_t$ se genera independientemente de las curvas, con probabilidades

\begin{equation}
\label{eq:ane_algB1_mixture}
P(Z_t=k)=\pi_k,
\qquad
\sum_{k=1}^{K}\pi_k=1.
\end{equation}

Por tanto, la heterogeneidad presente en las curvas proviene de la existencia de distintos mecanismos de generación, mientras que la asignación de cada período a un mecanismo no depende todavía de la curva anterior.\\

La generación procede de la siguiente forma:

\begin{enumerate}
    \item Se fijan los $K$ mecanismos, sus funciones medias $\mu_k(\tau)$ y sus coeficientes funcionales $\beta_{kj}(\tau)$.

    \item Se calculan los tres rasgos $R_{t-1,1}$, $R_{t-1,2}$ y $R_{t-1,3}$ a partir de la curva anterior.

    \item Se genera $Z_t$ de acuerdo con las probabilidades $\{\pi_k\}_{k=1}^{K}$.

    \item Se genera $X_t(\tau)$ mediante \eqref{eq:ane_algB1}, utilizando los parámetros correspondientes al mecanismo seleccionado.
\end{enumerate}

\begin{table}[H]
\centering
\caption{Parámetros del Algoritmo~B-1.}
\label{tab:ane_algB1}
\begin{tabular}{llc}
\hline
Símbolo & Descripción & Valor \\
\hline
$K$ & número de mecanismos funcionales & $3$ \\
$p$ & número de rasgos funcionales & $3$ \\
$\pi_k$ & probabilidad del mecanismo $k$ & $1/3$ \\
$w_j(\tau)$ & funciones de ponderación & prefijadas \\
$\beta_{kj}(\tau)$ & coeficientes funcionales & prefijados \\
\hline
\end{tabular}
\end{table}


\subsection{Algoritmo B-2: Mezcla de mecanismos funcionales no lineales}
\label{ane_00_02_02_alg_b2}

Este algoritmo extiende el escenario anterior incorporando mecanismos cuya relación entre los rasgos de la curva anterior y la curva futura puede ser no lineal. El objetivo es evaluar el comportamiento de los métodos cuando la heterogeneidad no se limita a diferencias en los parámetros de una relación lineal, sino que los distintos mecanismos presentan formas funcionales de dependencia diferentes.

La curva se genera mediante

\begin{equation}
\label{eq:ane_algB2}
X_t(\tau)
=
\mu_{Z_t}(\tau)
+
f_{Z_t}\!\left(
R_{t-1,1},R_{t-1,2},R_{t-1,3}
\right)(\tau)
+
\varepsilon_t(\tau),
\end{equation}

donde $f_k$ representa el mecanismo funcional asociado al estado $k$. Para introducir diferentes grados de no linealidad, se consideran mecanismos con estructuras distintas. Por ejemplo, un mecanismo puede presentar una relación lineal,

\begin{equation}
f_1(R_1,R_2,R_3)(\tau)
=
\beta_{11}(\tau)R_1
+
\beta_{12}(\tau)R_2
+
\beta_{13}(\tau)R_3,
\end{equation}

mientras que otros pueden incorporar términos cuadráticos o de interacción,

\begin{equation}
f_2(R_1,R_2,R_3)(\tau)
=
\beta_{21}(\tau)R_1^2
+
\beta_{22}(\tau)R_2
+
\beta_{23}(\tau)R_2R_3,
\end{equation}

y

\begin{equation}
f_3(R_1,R_2,R_3)(\tau)
=
\beta_{31}(\tau)R_1R_2
+
\beta_{32}(\tau)R_2^2
+
\beta_{33}(\tau)R_3^2.
\end{equation}

En este escenario, la asignación $Z_t$ mantiene probabilidades fijas, por lo que la complejidad adicional proviene exclusivamente de la forma de los mecanismos de generación.

La generación procede de forma análoga al Algoritmo~B-1, reemplazando la relación lineal por el mecanismo $f_{Z_t}$ correspondiente.

\begin{table}[H]
\centering
\caption{Parámetros del Algoritmo~B-2.}
\label{tab:ane_algB2}
\begin{tabular}{llc}
\hline
Símbolo & Descripción & Valor \\
\hline
$K$ & número de mecanismos funcionales & $3$ \\
$p$ & número de rasgos funcionales & $3$ \\
$\pi_k$ & probabilidad del mecanismo $k$ & $1/3$ \\
$f_1$ & mecanismo lineal & definido en el algoritmo \\
$f_2$ & mecanismo cuadrático & definido en el algoritmo \\
$f_3$ & mecanismo con interacciones & definido en el algoritmo \\
\hline
\end{tabular}
\end{table}


\subsection{Algoritmo B-3: Mezcla funcional con asignación dependiente de múltiples rasgos}
\label{ane_00_02_03_alg_b3}

Este algoritmo incorpora el nivel de complejidad más alto de los escenarios basados en procesos funcionales. Además de permitir mecanismos de generación no lineales, la selección del mecanismo deja de ser independiente de la curva anterior. En consecuencia, las características de $X_{t-1}$ determinan probabilísticamente qué mecanismo será utilizado para generar $X_t$.\\

Sea nuevamente $Z_t \in \{1,\ldots,K\}$ la variable latente. Sus probabilidades condicionales se definen mediante

\begin{equation}
\label{eq:ane_algB3_probs}
P(Z_t=k\mid X_{t-1})
=
\frac{
\exp\!\left\{
q_k(R_{t-1,1},R_{t-1,2},R_{t-1,3})
\right\}
}{
\displaystyle
\sum_{h=1}^{K}
\exp\!\left\{
q_h(R_{t-1,1},R_{t-1,2},R_{t-1,3})
\right\}
},
\end{equation}

donde $q_k(\cdot)$ determina cómo los distintos rasgos funcionales influyen en la probabilidad de seleccionar el mecanismo $k$.\\

Una vez seleccionado el mecanismo, la curva se genera mediante

\begin{equation}
\label{eq:ane_algB3}
X_t(\tau)
=
\mu_{Z_t}(\tau)
+
f_{Z_t}\!\left(
R_{t-1,1},R_{t-1,2},R_{t-1,3}
\right)(\tau)
+
\varepsilon_t(\tau).
\end{equation}

De este modo, la dependencia temporal presenta dos niveles. Primero, la curva anterior determina probabilísticamente el mecanismo que será utilizado. Segundo, el mecanismo seleccionado determina la forma en que los rasgos de la curva anterior se transforman en la nueva curva. Por tanto, la distribución condicional de $X_t$ es una mezcla cuyos pesos dependen de múltiples características funcionales de $X_{t-1}$.\\

La generación procede de la siguiente forma:

\begin{enumerate}
    \item Se fijan los mecanismos $f_1,\ldots,f_K$ y sus correspondientes funciones medias.

    \item Se calculan los rasgos funcionales $R_{t-1,1},R_{t-1,2},R_{t-1,3}$.

    \item Se calculan las probabilidades condicionales de los mecanismos mediante \eqref{eq:ane_algB3_probs}.

    \item Se genera $Z_t$ a partir de dichas probabilidades.

    \item Se genera $X_t(\tau)$ mediante \eqref{eq:ane_algB3}, utilizando el mecanismo seleccionado.
\end{enumerate}

\begin{table}[H]
\centering
\caption{Parámetros del Algoritmo~B-3.}
\label{tab:ane_algB3}
\begin{tabular}{llc}
\hline
Símbolo & Descripción & Valor \\
\hline
$K$ & número de mecanismos funcionales & $3$ \\
$p$ & número de rasgos funcionales & $3$ \\
$q_k$ & función de asignación del mecanismo $k$ & prefijada \\
$f_k$ & mecanismo funcional no lineal & prefijado \\
\hline
\end{tabular}
\end{table}

\section{Métodos basados en Coeficientes de la Representación}
\label{ane_00_03_metodos_coeficientes_bases}
Este tercer enfoque simula directamente los coeficientes de la representación de las curvas. Sea $\{\varphi_j\}_{j=1}^{J}$ un sistema ortonormal fijo en $L^2(\mathcal{T})$, en este caso la base de Fourier, y $\mu(\tau)$ la función media. Cada curva se escribe como
\begin{equation}\label{eq:ane_coef_curva}
    X_t(\tau) = \mu(\tau) + \sum_{j=1}^{J} a_{tj}\,\varphi_j(\tau),
    \qquad
    a_{tj} = \sqrt{\lambda_j}\;\tilde a_{tj},
\end{equation}
donde $a_{tj}$ es el coeficiente de la $j$-ésima componente, $\lambda_1 > \dots > \lambda_J > 0$ es un espectro prefijado que fija la varianza marginal de cada componente, y $\tilde a_{tj}$ es una serie de tiempo de media cero y varianza unitaria. Toda la dependencia temporal del proceso, lineal y no lineal, se especifica sobre las series $\{\tilde a_{tj}\}$.\\

Como el sistema $\{\varphi_j\}$ es fijo y conocido, los coeficientes simulados se entregan directamente a los métodos, en el papel de los scores de la representación, sin construir las curvas ni aplicar el esquema de observación. De este modo se aísla el efecto de la estructura de dependencia y del número de scores retenidos, sin confundirlo con el error de estimar la representación ni con el ruido de medición. \\


Los algoritmos son los siguientes:
\begin{itemize}
    \item El Algoritmo~C-1  tiene una dependencia lineal débil y dependencia no lineal entre las dos primeras componentes.
    \item El Algoritmo~C-2 refuerza la dependencia lineal, dejando intacta la no lineal.
    \item El Algoritmo~C-3 traslada la dependencia a componentes subordinadas.
\end{itemize}

\subsection{Algoritmo C-1: Proceso con dependencia no lineal entre coeficientes}
\label{ane_00_03_01_alg_c1}
Este algoritmo genera coeficientes cuya dependencia temporal es débil en el sentido lineal, pero fuerte en el no lineal, con el fin de evaluar la capacidad de los métodos para aprovechar dependencia que un modelo lineal no detecta. Sea $(j^\ast, j^\ast+1)$ el par de componentes que interactúan, y sean $\{e_{tj}\}$ variables gaussianas estándar independientes. La dinámica de las series $\tilde a_{tj}$ es
\begin{align}
    \tilde a_{t,j^\ast} &= \phi\, \tilde a_{t-1,j^\ast} + \sqrt{1-\phi^2}\; e_{t,j^\ast}, \label{eq:ane_algC1_lineal}\\
    \tilde a_{t,j^\ast+1} &= c\, \tilde a_{t-1,j^\ast} + b\,\frac{\tilde a_{t-1,j^\ast}^{2}-1}{\sqrt{2}} + \sqrt{1-c^2-b^2}\; e_{t,j^\ast+1}, \label{eq:ane_algC1_cuadratica}
\end{align}
y $\tilde a_{tj} = e_{tj}$ para el resto de las componentes. La primera ecuación es un autorregresivo lineal de coeficiente $\phi$. En la segunda, $c$ es el coeficiente de la dependencia lineal cruzada y $b$ el de la dependencia cuadrática. Las varianzas de innovación se despejan para que cada $\tilde a_{tj}$ tenga varianza unitaria y el espectro marginal quede fijado en $\lambda_j$. Un predictor de $\tilde a_{t,j^\ast+1}$ lineal en $\tilde a_{t-1,j^\ast}$ explica una fracción $c^2$ de su varianza, mientras que uno que capture la relación cuadrática explica $c^2 + b^2$. Las demás autocorrelaciones lineales de rezago~1 entre componentes son de orden menor.\\

La generación procede de la siguiente forma:
\begin{enumerate}
    \item Se fijan el espectro $\{\lambda_j\}_{j=1}^{J}$ y los parámetros $j^\ast$, $\phi$, $c$ y $b$.

    \item Se extrae $\tilde a_{0,j^\ast} \sim \mathcal{N}(0,1)$ y se itera \eqref{eq:ane_algC1_lineal}--\eqref{eq:ane_algC1_cuadratica} durante $T$ períodos, obteniendo $\tilde{\mathbf{A}} \in \mathbb{R}^{T \times J}$, y se escalan las columnas para obtener $\mathbf{A} = \tilde{\mathbf{A}}\,\mathrm{diag}(\sqrt{\lambda_1}, \dots, \sqrt{\lambda_J})$.

    \item La matriz de coeficientes $\mathbf{A}$ es el producto del algoritmo y se entrega directamente a los métodos.
\end{enumerate}

\begin{table}[H]
\centering
\caption{Parámetros del Algoritmo~C-1.}
\label{tab:ane_algC1}
\begin{tabular}{llc}
\hline
Símbolo & Descripción & Valor \\
\hline
$J$             & términos de la base de Fourier                            & $10$ \\
$\rho$          & razón del espectro geométrico $\lambda_j=\rho^{\,j-1}$    & $0.5$ \\
$j^\ast$        & primera componente del par que interactúa                 & $1$ \\
$\phi$          & coeficiente autorregresivo de $\tilde a_{t,j^\ast}$       & $0.2$ \\
$c$             & coeficiente de la dependencia lineal cruzada              & $0.15$ \\
$b$             & coeficiente de la dependencia cuadrática                  & $0.5$ \\
\hline
\end{tabular}
\end{table}

\subsection{Algoritmo C-2: Proceso con dependencia lineal reforzada}
\label{ane_00_03_02_alg_c2}
Este algoritmo genera coeficientes con la misma construcción del Algoritmo~C-1, aumentando el coeficiente autorregresivo $\phi$ y dejando intacta la dependencia cuadrática, con el fin de aislar cuánto del desempeño de cada método proviene de la dependencia lineal y cuánto de la no lineal. La generación sigue los pasos del Algoritmo~C-1.

\begin{table}[H]
\centering
\caption{Parámetros del Algoritmo~C-2. Coinciden con los del Cuadro~\ref{tab:ane_algC1}, salvo $\phi$.}
\label{tab:ane_algC2}
\begin{tabular}{llc}
\hline
Símbolo & Descripción & Valor \\
\hline
$\phi$  & coeficiente autorregresivo de $\tilde a_{t,j^\ast}$ & $0.45$ \\
\hline
\end{tabular}
\end{table}

\subsection{Algoritmo C-3: Proceso con dependencia en componentes subordinadas}
\label{ane_00_03_03_alg_c3}
Este algoritmo genera coeficientes con la misma construcción del Algoritmo~C-1, trasladando el par de componentes que interactúan a las componentes 3 y 4, con el fin de evaluar el desempeño de los métodos cuando la información sobre el futuro reside en componentes de menor varianza. Con el espectro geométrico empleado, las dos primeras componentes concentran cerca del $75\,\text{\%}$ de la varianza y estas dos, cerca del $19\,\text{\%}$, por lo que un método que retenga pocos scores no accede a la dependencia. La generación sigue los pasos del Algoritmo~C-1.

\begin{table}[H]
\centering
\caption{Parámetros del Algoritmo~C-3. Coinciden con los del Cuadro~\ref{tab:ane_algC1}, salvo $j^\ast$.}
\label{tab:ane_algC3}
\begin{tabular}{llc}
\hline
Símbolo & Descripción & Valor \\
\hline
$j^\ast$    & primera componente del par que interactúa   & $3$ \\
\hline
\end{tabular}
\end{table}